\section{Empirical Strategy}
\label{sec:empirical}

\subsection{OLS Baseline (H1)}
\label{sec:empirical_ols}

Our baseline specification estimates the association between FL presence
and tender outcomes:
\begin{equation}
  y_{igt} = \beta \cdot \text{losers}_{igt} + \mathbf{x}_{igt}'
  \boldsymbol{\delta} + \alpha_g + \lambda_t + \gamma_k + \varepsilon_{igt}
  \label{eq:main}
\end{equation}
where $y_{igt}$ is the outcome for tender-item $i$ in item group $g$ at
time $t$ and purchasing unit $k$; $\text{losers}_{igt} \in \{0,1\}$
indicates FL presence; $\mathbf{x}_{igt}$ includes a convite indicator;
$\alpha_g$, $\lambda_t$, and $\gamma_k$ are item, year, and PBU fixed
effects; and $\varepsilon_{igt}$ is the error term clustered at the item
level.

We estimate four dependent variables: (i) $\log$ negotiated price, (ii)
$\log$ number of firms, (iii) $\log$ number of bids, and (iv) $\log$
number of non-FL (``genuine'') firms. For each DV we report four
specifications: (1) item + year FE, (2) item + year + PBU FE, (3)
preg\~ao only with all FE, (4) convite only with all FE.

The OLS coefficient $\hat{\beta}$ identifies a conditional
association---the average difference in outcomes between FL-present and
FL-absent tenders within the same item type, year, and (in
specifications 2--4) purchasing unit. Under H1, $\hat{\beta} > 0$ for
prices establishes the FL screen's empirical validity without requiring a
causal interpretation.

\subsection{Instrumental Variable Strategy (H2)}
\label{sec:empirical_iv}

To move beyond association toward causation (H2), we instrument FL
presence with a \emph{leave-one-out} measure of FL supply:
\begin{equation}
  Z_{kgt} = \sum_{j \neq k} \mathbf{1}[\text{FL firm active at PBU } j
  \text{ in group } g, \text{ year } t]
  \label{eq:iv}
\end{equation}
For each tender-item at PBU $k$ in item group $g$ and year $t$, the
instrument counts FL firms that are active at \emph{other} PBUs in the
same product market and year. The identifying assumption is that FL
activity at distant PBUs affects PBU $k$'s outcomes only through the
probability of FL participation at $k$ itself---a supply-side exclusion
restriction in the spirit of \citet{goldsmithpinkham2020bartik}.

The first stage is:
\begin{equation}
  \text{losers}_{igt} = \pi \cdot Z_{kgt} + \mathbf{x}_{igt}'
  \boldsymbol{\phi} + \tilde{\alpha}_g + \tilde{\lambda}_t +
  \tilde{\gamma}_k + u_{igt}
  \label{eq:first_stage}
\end{equation}
We report the first-stage $F$-statistic and partial $R^2$ to assess
instrument strength. We also report a \emph{placebo IV} using
sub-threshold always-losers (firms that never win but participate below
the IQR threshold) as the instrument source; under H2, these firms
should not drive the same price effects.

\paragraph{Threats to the exclusion restriction.}
The identifying assumption requires that FL activity at other PBUs
affects PBU $k$'s outcomes \emph{only} through FL participation at $k$.
Three threats merit discussion. First, \emph{correlated demand shocks}:
if the same item-group experiences price shocks statewide, FL supply at
other PBUs may correlate with prices at PBU $k$ directly. The item,
year, and PBU fixed effects absorb time-invariant item characteristics,
aggregate trends, and PBU-specific price levels. As a further test, we
estimate the 2SLS specification with item-group $\times$ year fixed
effects, which absorb product-market-year shocks and restrict
identification to within-cell variation across PBUs
(Table~\ref{tab:iv_main}, Panel~C). Second, \emph{statewide cartel
organization}: if a cartel coordinates FL deployment across PBUs
simultaneously, the LOO instrument captures coordinated entry rather than
independent supply variation. We address this by testing balance:
Table~\ref{tab:iv_balance} regresses pre-determined observables on the
instrument and finds no significant associations, consistent with
quasi-random FL supply. Third, \emph{enforcement variation}: if
enforcement intensity varies across PBUs and correlates with both FL
supply and prices, the instrument would be invalid. The PBU fixed effects
absorb time-invariant enforcement differences, and the item-group
$\times$ year specification further absorbs any enforcement changes that
are product-market specific.

\subsection{Network-Based FL Classification}
\label{sec:empirical_network}

Not all 2,735 FL firms need be cover bidders. Some may be genuinely
incompetent firms that persistently bid too high. To explore
heterogeneity, we analyze co-bidding networks using the firm--tender map
(16.8 million rows):

\begin{enumerate}
  \item For each FL firm, we identify all non-FL co-bidders and
    co-winners across all tenders.
  \item We compute four metrics:
    \begin{itemize}
      \item \emph{Winner HHI}: concentration of firms that win tenders
        where this FL firm bids. High HHI indicates the FL firm operates
        in concentrated markets with few distinct winners.
      \item \emph{Repeat partners}: number of non-FL firms that co-bid
        with this FL firm in $\geq 3$ tenders.
      \item \emph{Top partner share}: fraction of the FL firm's tenders
        shared with its most frequent co-bidder.
      \item \emph{Bid spread CV}: coefficient of variation of the FL
        firm's bid-to-winning-bid ratio.
    \end{itemize}
  \item We classify FL firms as \emph{high-concentration} if winner HHI
    exceeds the median (0.133) and repeat partners $\geq 2$, yielding
    1,356 high-concentration and 1,379 low-concentration FL firms.
\end{enumerate}

We then split the treatment into two
indicators---$\text{has\_high\_susp\_fl}$ and
$\text{has\_low\_susp\_fl}$---and test whether the price effect varies
by market structure.

\subsection{Bajari--Ye Test Design}
\label{sec:empirical_bajari}

We implement the \citet{bajari2003deciding} tests for bid coordination.
The first stage regresses $\log$ bid values on observable cost shifters
(reference price, firm size, item and year fixed effects) to obtain bid
residuals. We report the first-stage $R^2$ and coefficient estimates in
Appendix~\ref{app:robustness}.

Using these residuals, we conduct three tests:

\begin{enumerate}
  \item \textbf{Exchangeability}: Kolmogorov--Smirnov test comparing the
    distributions of FL and non-FL bid residuals. Under competitive
    bidding with symmetric information, all bidders' residuals should
    come from the same distribution.
  \item \textbf{Conditional independence}: for each tender with $\geq 2$
    FL bids, compute the mean product of pairwise FL residuals. Under
    independence, this product should be zero in expectation. We compare
    the FL pairwise product with the non-FL pairwise product using a
    tender-level bootstrap (1,000 iterations).
  \item \textbf{Fake-groups placebo}: randomly assign non-FL firms into
    two groups within each tender and repeat both tests. If the
    methodology is valid, the placebo should yield null results.
\end{enumerate}

The pairwise product statistic has a natural interpretation: a value of
$\bar{r}_i \cdot \bar{r}_j = 5.16$ means that when one FL bid is R\$1
above its predicted value, the paired FL bid is on average R\$5.16 above
its prediction, indicating systematic co-movement.
